% from ega4-18-fr.tex
\oldpage[IV-4]{187}

de $X$ (ce qui entraîne, par (19.1.1), que $f$ est une immersion régulière au point $g(z)$); alors, pour que $f\circ g$ soit une immersion régulière au point $z$, il faut et il suffit que $g$ soit une immersion régulière au point $z$.
\end{enumerate}
\end{proposition}

L'assertion (i) est une autre manière d'exprimer (16.1.10). Pour démontrer la première assertion de (ii), on peut se borner au cas où $Y$ est un sous-préschéma fermé de $X$ défini par un Idéal régulier (resp. quasi-régulier) $\mathscr J$; alors $Y'$ est le sous-préschéma fermé de $X'$ défini par l'Idéal $g^*(\mathscr J)\sh{O}_{X'}$ (I, 4.4.5) qui s'identifie ici à $g^*(\mathscr J)$ puisque $g$ est plat; en remplaçant $X$ par un ouvert affine assez petit, on peut supposer que $\mathscr J$ admet une suite régulière (resp. quasi-régulière) de générateurs $f_i$ (sections de $\sh{O}_X$ au-dessus de $X$). Si la suite $(f_i)$ est régulière, il en est de même de la suite $(f_i\otimes1)$ (qui engendre $g^*(\mathscr J)\sh{O}_{X'}$) en vertu de ($0$, 15.2.5); le résultat analogue pour les suites quasi-régulières découle aussitôt de la platitude de $g$ et du critère (16.9.4). De même, si $g$ est fidèlement plat et quasi-compact, et si $g^*(\mathscr J)$ est quasi-régulier, $\mathscr J$ est de type fini en vertu de (2.5.2); par platitude, on a $g^*(\mathscr J)/g^*(\mathscr J)^2=g^*(\mathscr J/\mathscr J^2)$, donc, puisque $g^*(\mathscr J/\mathscr J^2)$ est localement libre par hypothèse (16.9.4), il en est de même de $\mathscr J/\mathscr J^2$ (2.5.2). Enfin, la condition (iii) de (16.9.4) relative à $g^*(\mathscr J)$ entraîne la même condition pour $\mathscr J$ en vertu de la fidèle platitude de $g$ (2.2.7). On conclut donc de (16.9.4) que $\mathscr J$ est quasi-régulier.

Pour prouver (iii), on peut de même supposer que $Y$ et $Z$ sont des sous-préschémas fermés de $X$, définis respectivement par des Idéaux $\mathscr J$, $\mathscr K$ de $\sh{O}_X$ tels que $\mathscr J\subset\mathscr K$; de plus, on peut supposer qu'il existe des sections $f_i$ ($1\leq i\leq m$) de $\mathscr J$ au-dessus de $X$ engendrant $\mathscr J$ et telles que pour $1\leq i\leq m$, l'endomorphisme de $\sh{O}_X/(\sum_{j=1}^{i-1}f_j\sh{O}_X)$ défini par $f_i$ soit injectif, et des sections $\overline f_i$ ($m+1\leq i\leq n$) de $\mathscr K/\mathscr J$ au-dessus de $X$ engendrant $\mathscr K/\mathscr J$ et telles que pour $m+1\leq i\leq n$, l'endomorphisme de
\[
  (\sh{O}_X/\mathscr J)\big/\Bigl(\sum_{j=m+1}^{i-1}\overline f_j(\sh{O}_X/\mathscr J)\Bigr)
\]
défini par $\overline f_i$ soit injectif. Pour tout $z\in Z$, soit $U$ un voisinage ouvert de $z$ dans $X$ tel qu'il existe une section $f_i$ de $\mathscr K$ au-dessus de $U$, dont $\overline f_i|U$ soit la classe dans $\Gamma(U,\mathscr K/\mathscr J)$ pour $m+1\leq i\leq n$; alors
\[
  (\sh{O}_X/\mathscr J)\big/\Bigl(\sum_{j=m+1}^{i-1}\overline f_j(\sh{O}_X/\mathscr J)\Bigr),
\]
restreint à $U$ est isomorphe à
\[
  \sh{O}_U\big/\Bigl((\mathscr J|U)+\sum_{j=m+1}^{i-1}f_j\sh{O}_U\Bigr)
  =\sh{O}_U\big/\Bigl(\sum_{j=1}^{i-1}f_j\sh{O}_U\Bigr),
\]
et on voit donc que la suite $(f_i)_{1\leq i\leq n}$ est régulière dans $\sh{O}_U$; comme elle engendre $\mathscr K|U$, cela prouve la première assertion de (iii). Pour prouver la dernière assertion de (iii), on peut se borner (avec les mêmes notations) au cas où les images $t_i$ des $f_i$ dans $\mathscr K/\mathscr K^2$ (pour $1\leq i\leq n$) forment une base de cet $(\sh{O}_X/\mathscr K)$-Module (16.9.5); pour la même raison, on peut supposer que les $t_i$ tels que $1\leq i\leq m$ sont les images canoniques des éléments d'une base du $(\sh{O}_X/\mathscr K)$-Module $g^*(\mathscr J/\mathscr J^2)$, et que les images canoniques $\overline t_i$ des $t_i$ dans $(\mathscr K/\mathscr J)/(\mathscr K/\mathscr J)^2$ pour $m+1\leq i\leq n$ forment une base de cet $(\sh{O}_X/\mathscr K)$-Module. Cela achève de prouver (iii).

Passons à la démonstration de (iv). Le fait que $a)$ implique $c)$ résulte de (iii). Prouvons que $c)$ implique $a)$. Posons encore $A=\sh{O}_{X,f(g(z))}$, de sorte que $\sh{O}_{Y,g(z)}=A/\mathfrak J$,

